% !TITLE 桂枝香·金陵怀古
% !AUTHOR 王安石
% !DYNASTY 两宋
% !GENRE 词
% !ORDER 9
% !EDITION 《全宋词》，唐圭璋编，中华书局
% !EDITION_STATUS pending

\begin{poem}
登临送目\textsuperscript{1}，正故国\textsuperscript{2}晚秋，天气初肃\textsuperscript{3}。
千里澄江似练\textsuperscript{4}，翠峰如簇\textsuperscript{5}。
归帆去棹残阳里，背西风、酒旗斜矗\textsuperscript{6}。
彩舟云淡，星河鹭起\textsuperscript{7}，画图难足。

念往昔、繁华竞逐。叹门外楼头\textsuperscript{8}，悲恨相续。
千古凭高对此，谩嗟荣辱\textsuperscript{9}。
六朝\textsuperscript{10}旧事随流水，但寒烟、芳草凝绿。
至今商女\textsuperscript{11}，时时犹唱，后庭遗曲\textsuperscript{12}。
\end{poem}

\begin{annotations}
\notetext{1}{送目：放眼远望。登高远眺金陵（今南京）的秋景。}
\notetext{2}{故国：旧时的都城。金陵是六朝古都，所以叫“故国”。}
\notetext{3}{初肃：刚开始变得萧索。肃，肃杀、冷落。}
\notetext{4}{澄江似练：化用谢朓名句“澄江静如练”。清澈的长江像一条白练。}
\notetext{5}{翠峰如簇：翠绿的山峰像箭头一样聚集在一起。簇（cù），簇拥聚集。}
\notetext{6}{酒旗斜矗：酒家的旗幡斜斜地竖着。矗（chù），直立。}
\notetext{7}{星河鹭起：银河（指长江）上有白鹭飞起。}
\notetext{8}{门外楼头：化用杜牧《台城曲》“门外韩擒虎，楼头张丽华”。隋将韩擒虎已到门外，陈后主还在楼头与宠妃张丽华寻欢——写亡国的荒诞与悲恨。}
\notetext{9}{谩嗟荣辱：徒然地感叹兴衰荣辱。谩（màn），徒然、白白地。}
\notetext{10}{六朝：东吴、东晋、宋、齐、梁、陈，六个建都金陵的朝代。六朝的繁华都已随流水而去。}
\notetext{11}{商女：酒家卖唱的歌女。}
\notetext{12}{后庭遗曲：指陈后主所作的《玉树后庭花》，被后世视为“亡国之音”。化用杜牧《泊秦淮》“商女不知亡国恨，隔江犹唱后庭花”。}
\end{annotations}

\begin{appreciation}
《桂枝香·金陵怀古》是王安石晚年罢相、退居金陵时所作，历来被推为金陵怀古词的压卷之作之一。据说苏轼读到这首词，也忍不住叹服，说王安石是“野狐精”——两人政见水火不容而能让对手说出这句话可见其分量。

上阕写金陵秋色。“登临送目，正故国晚秋，天气初肃。”登高放眼，正赶上旧都的深秋，天气刚刚显出萧索。“故国”指金陵，东吴、东晋、宋、齐、梁、陈六朝都建都于此，所以称旧都。“千里澄江似练，翠峰如簇”，江水清澈如一条白练铺向千里，青翠的山峰像箭镞一样聚拢。这一句的来历很讲究：南朝谢朓写过“澄江静如练”，王安石几乎是原样取来，站在谢朓望过的江边，用的还是谢朓的句子。“归帆去棹残阳里，背西风、酒旗斜矗”，斜阳里帆影来去，酒旗背着西风斜斜地挑着。“彩舟云淡，星河鹭起，画图难足”，白鹭从江面飞起，仿佛银河里溅起的星点。他说这样的景致再好的画笔也描摹不尽。

下阕陡然一转。“念往昔、繁华竞逐。”金陵六朝，一个比一个繁华，又一个接一个覆亡。“叹门外楼头，悲恨相续。”这一句用杜牧的典故：隋将韩擒虎已经打到了宫门外，陈后主还在楼头与宠妃张丽华听歌——亡国竟亡得这样荒唐。“千古凭高对此，谩嗟荣辱”，千年来登高望江的人也只能白白感叹一番兴亡。“六朝旧事随流水，但寒烟、芳草凝绿”，六朝的繁华都随江水流走，剩下的只是一片寒烟、满目芳草。末句更重：“至今商女，时时犹唱，后庭遗曲。”《玉树后庭花》是陈后主所作，后世把它当作亡国之音，前面读过杜牧的《泊秦淮》，王安石在这里又听见了一遍。六朝早已亡了，歌却还在唱。

从谢朓到杜牧，再到王安石，一代代人站在金陵看江水，看王朝起落。而王安石写这首词的时候，北宋自己也正走下坡路：他的变法失败了而他看得清清楚楚却什么也拦不住。这才是怀古最深的悲凉——历史从来不听劝。
\end{appreciation}
